%% kapa: exam=1 round=28 year=2017 period=2 form=A subject=appraisal_law subject_ko=감정평가관계법규 seq=37 number=37
%% source: Q-Net 감정평가사 기출문제 (cst00309) artlSeq=5199603 file=extracted/1008/5199603/2017년 제28회 감정평가사 1차시험 문제지/2교시 A형.pdf
%% conversion: kordoc 4.24.51 -> tools/gen1.py (auto); review=proofread-sonnet; flags=none
\begin{question}{37}
도시 및 주거환경정비법령상 조합을 설립하는 경우 토지등소유자의 동의자수 산정방법으로 옳지 않은 것은?
\begin{choices}
  \item 주택재건축사업의 경우 1명이 둘 이상의 소유권을 소유하고 있는 경우에는 소유권의 수에 관계없이 토지등소유자를 1명으로 산정한다.
  \item 주택재개발사업의 경우 하나의 건축물이 수인의 공유에 속하는 때에는 그 수인을 대표하는 1인을 토지등소유자로 산정한다.
  \item 국ㆍ공유지에 대해서는 그 재산관리청을 토지등소유자로 산정한다.
  \item 주택재개발사업의 경우 토지에 지상권이 설정되어 있는 경우에는 토지의 소유자와 해당 토지의 지상권자를 대표하는 1인을 토지등소유자로 산정한다.
  \item 도시환경정비사업의 경우 토지등소유자가 정비구역 지정 후에 정비사업을 목적으로 토지를 추가로 취득하여 1인이 다수 필지의 토지를 소유하게 된 경우에는 필지의 수에 관계없이 토지등소유자를 1인으로 산정한다.
\end{choices}
\end{question}
\answer{5}
